%% figstyle.tex
%% 卒論用 pgfplots 図の共通設定。
%% 本文 (1029307812_特別研究報告書2026.tex) のプリアンブルで
%%   %% figstyle.tex
%% 卒論用 pgfplots 図の共通設定。
%% 本文 (1029307812_特別研究報告書2026.tex) のプリアンブルで
%%   %% figstyle.tex
%% 卒論用 pgfplots 図の共通設定。
%% 本文 (1029307812_特別研究報告書2026.tex) のプリアンブルで
%%   %% figstyle.tex
%% 卒論用 pgfplots 図の共通設定。
%% 本文 (1029307812_特別研究報告書2026.tex) のプリアンブルで
%%   \input{analysis/figstyle}
%% を呼び、各図は \input{analysis/fig_xxx} で取り込む。
%%
%% ※ documentclass のオプションに [dvipdfmx] が付いているので、
%%   tikz/pgfplots は自動で dvipdfmx ドライバを使う（platex + dvipdfmx 構成）。

\usepackage{tikz}
\usepackage{pgfplots}
\usepackage{pgfplotstable}
\pgfplotsset{compat=1.18}
\usepgfplotslibrary{groupplots}
\usetikzlibrary{pgfplots.groupplots}

%% データファイルの場所（本文からの相対パス）。必要なら上書き可。
\providecommand{\figdata}{analysis/data}

%% 図が多い節で float が後ろの節へ流れ込まないように配置を緩める。
\renewcommand{\topfraction}{0.9}
\renewcommand{\bottomfraction}{0.9}
\renewcommand{\textfraction}{0.08}
\renewcommand{\floatpagefraction}{0.75}
\setcounter{topnumber}{2}
\setcounter{bottomnumber}{1}
\setcounter{totalnumber}{4}

%% '#' 始まりの行をコメントとして無視させる（gnuplot 風 .dat 用）。
\pgfplotstableset{col sep=space, comment chars={\#}}

%% beta = 0.5,1,2,3 を表す色（色覚多様性に配慮した4色）
\definecolor{betaA}{HTML}{0072B2} % beta=0.5  青
\definecolor{betaB}{HTML}{009E73} % beta=1    緑
\definecolor{betaC}{HTML}{E69F00} % beta=2    橙
\definecolor{betaD}{HTML}{D55E00} % beta=3    朱

%% 図の基本スタイル
\pgfplotsset{
  thesisaxis/.style={
    width=7.2cm, height=5.6cm,
    scale only axis,
    tick align=outside,
    tick label style={font=\footnotesize},
    label style={font=\small},
    legend style={font=\footnotesize, draw=none, fill opacity=0.85,
      text opacity=1, row sep=1pt, inner sep=2pt},
    every axis plot/.append style={line width=0.9pt, mark size=1.6pt},
    grid=both,
    grid style={line width=0.2pt, draw=gray!18},
    major grid style={line width=0.3pt, draw=gray!30},
  },
  % beta 4本のカーブを一括設定する系列スタイル
  betaplotA/.style={color=betaA, mark=*,        mark options={fill=betaA}},
  betaplotB/.style={color=betaB, mark=square*,  mark options={fill=betaB}},
  betaplotC/.style={color=betaC, mark=triangle*,mark options={fill=betaC}},
  betaplotD/.style={color=betaD, mark=diamond*, mark options={fill=betaD}},
}

%% delta 軸（対数）の共通設定
\pgfplotsset{
  deltaxaxis/.style={
    xmode=log, log basis x=10,
    xtick={0.3333,1,2,3,5,10},
    xticklabels={$\tfrac13$,$1$,$2$,$3$,$5$,$10$},
    xlabel={$\frac{|\Delta \alpha|E}{p}$},
    xmin=0.28, xmax=12,
  },
}

%% beta 凡例（4系列に共通で使う）
\newcommand{\betalegend}{\legend{$\beta pE=0.5$,$\beta pE=1$,$\beta pE=2$,$\beta pE=3$}}

%% 4本の beta カーブ（誤差棒つき）を一括描画。
%%   #1 データファイル, #2 x列名, #3 値の接頭辞, #4 誤差の接頭辞
\newcommand{\plotbetaserr}[4]{%
  \addplot[betaplotA, error bars/.cd, y dir=both, y explicit]
    table[x=#2, y=#3_b05, y error=#4_b05]{#1};
  \addplot[betaplotB, error bars/.cd, y dir=both, y explicit]
    table[x=#2, y=#3_b1,  y error=#4_b1]{#1};
  \addplot[betaplotC, error bars/.cd, y dir=both, y explicit]
    table[x=#2, y=#3_b2,  y error=#4_b2]{#1};
  \addplot[betaplotD, error bars/.cd, y dir=both, y explicit]
    table[x=#2, y=#3_b3,  y error=#4_b3]{#1};
}

%% 4本の beta カーブ（誤差棒なし）。#1 ファイル, #2 x列, #3 値の接頭辞
\newcommand{\plotbetas}[3]{%
  \addplot[betaplotA] table[x=#2, y=#3_b05]{#1};
  \addplot[betaplotB] table[x=#2, y=#3_b1]{#1};
  \addplot[betaplotC] table[x=#2, y=#3_b2]{#1};
  \addplot[betaplotD] table[x=#2, y=#3_b3]{#1};
}

%% を呼び、各図は \input{analysis/fig_xxx} で取り込む。
%%
%% ※ documentclass のオプションに [dvipdfmx] が付いているので、
%%   tikz/pgfplots は自動で dvipdfmx ドライバを使う（platex + dvipdfmx 構成）。

\usepackage{tikz}
\usepackage{pgfplots}
\usepackage{pgfplotstable}
\pgfplotsset{compat=1.18}
\usepgfplotslibrary{groupplots}
\usetikzlibrary{pgfplots.groupplots}

%% データファイルの場所（本文からの相対パス）。必要なら上書き可。
\providecommand{\figdata}{analysis/data}

%% 図が多い節で float が後ろの節へ流れ込まないように配置を緩める。
\renewcommand{\topfraction}{0.9}
\renewcommand{\bottomfraction}{0.9}
\renewcommand{\textfraction}{0.08}
\renewcommand{\floatpagefraction}{0.75}
\setcounter{topnumber}{2}
\setcounter{bottomnumber}{1}
\setcounter{totalnumber}{4}

%% '#' 始まりの行をコメントとして無視させる（gnuplot 風 .dat 用）。
\pgfplotstableset{col sep=space, comment chars={\#}}

%% beta = 0.5,1,2,3 を表す色（色覚多様性に配慮した4色）
\definecolor{betaA}{HTML}{0072B2} % beta=0.5  青
\definecolor{betaB}{HTML}{009E73} % beta=1    緑
\definecolor{betaC}{HTML}{E69F00} % beta=2    橙
\definecolor{betaD}{HTML}{D55E00} % beta=3    朱

%% 図の基本スタイル
\pgfplotsset{
  thesisaxis/.style={
    width=7.2cm, height=5.6cm,
    scale only axis,
    tick align=outside,
    tick label style={font=\footnotesize},
    label style={font=\small},
    legend style={font=\footnotesize, draw=none, fill opacity=0.85,
      text opacity=1, row sep=1pt, inner sep=2pt},
    every axis plot/.append style={line width=0.9pt, mark size=1.6pt},
    grid=both,
    grid style={line width=0.2pt, draw=gray!18},
    major grid style={line width=0.3pt, draw=gray!30},
  },
  % beta 4本のカーブを一括設定する系列スタイル
  betaplotA/.style={color=betaA, mark=*,        mark options={fill=betaA}},
  betaplotB/.style={color=betaB, mark=square*,  mark options={fill=betaB}},
  betaplotC/.style={color=betaC, mark=triangle*,mark options={fill=betaC}},
  betaplotD/.style={color=betaD, mark=diamond*, mark options={fill=betaD}},
}

%% delta 軸（対数）の共通設定
\pgfplotsset{
  deltaxaxis/.style={
    xmode=log, log basis x=10,
    xtick={0.3333,1,2,3,5,10},
    xticklabels={$\tfrac13$,$1$,$2$,$3$,$5$,$10$},
    xlabel={$\frac{|\Delta \alpha|E}{p}$},
    xmin=0.28, xmax=12,
  },
}

%% beta 凡例（4系列に共通で使う）
\newcommand{\betalegend}{\legend{$\beta pE=0.5$,$\beta pE=1$,$\beta pE=2$,$\beta pE=3$}}

%% 4本の beta カーブ（誤差棒つき）を一括描画。
%%   #1 データファイル, #2 x列名, #3 値の接頭辞, #4 誤差の接頭辞
\newcommand{\plotbetaserr}[4]{%
  \addplot[betaplotA, error bars/.cd, y dir=both, y explicit]
    table[x=#2, y=#3_b05, y error=#4_b05]{#1};
  \addplot[betaplotB, error bars/.cd, y dir=both, y explicit]
    table[x=#2, y=#3_b1,  y error=#4_b1]{#1};
  \addplot[betaplotC, error bars/.cd, y dir=both, y explicit]
    table[x=#2, y=#3_b2,  y error=#4_b2]{#1};
  \addplot[betaplotD, error bars/.cd, y dir=both, y explicit]
    table[x=#2, y=#3_b3,  y error=#4_b3]{#1};
}

%% 4本の beta カーブ（誤差棒なし）。#1 ファイル, #2 x列, #3 値の接頭辞
\newcommand{\plotbetas}[3]{%
  \addplot[betaplotA] table[x=#2, y=#3_b05]{#1};
  \addplot[betaplotB] table[x=#2, y=#3_b1]{#1};
  \addplot[betaplotC] table[x=#2, y=#3_b2]{#1};
  \addplot[betaplotD] table[x=#2, y=#3_b3]{#1};
}

%% を呼び、各図は \input{analysis/fig_xxx} で取り込む。
%%
%% ※ documentclass のオプションに [dvipdfmx] が付いているので、
%%   tikz/pgfplots は自動で dvipdfmx ドライバを使う（platex + dvipdfmx 構成）。

\usepackage{tikz}
\usepackage{pgfplots}
\usepackage{pgfplotstable}
\pgfplotsset{compat=1.18}
\usepgfplotslibrary{groupplots}
\usetikzlibrary{pgfplots.groupplots}

%% データファイルの場所（本文からの相対パス）。必要なら上書き可。
\providecommand{\figdata}{analysis/data}

%% 図が多い節で float が後ろの節へ流れ込まないように配置を緩める。
\renewcommand{\topfraction}{0.9}
\renewcommand{\bottomfraction}{0.9}
\renewcommand{\textfraction}{0.08}
\renewcommand{\floatpagefraction}{0.75}
\setcounter{topnumber}{2}
\setcounter{bottomnumber}{1}
\setcounter{totalnumber}{4}

%% '#' 始まりの行をコメントとして無視させる（gnuplot 風 .dat 用）。
\pgfplotstableset{col sep=space, comment chars={\#}}

%% beta = 0.5,1,2,3 を表す色（色覚多様性に配慮した4色）
\definecolor{betaA}{HTML}{0072B2} % beta=0.5  青
\definecolor{betaB}{HTML}{009E73} % beta=1    緑
\definecolor{betaC}{HTML}{E69F00} % beta=2    橙
\definecolor{betaD}{HTML}{D55E00} % beta=3    朱

%% 図の基本スタイル
\pgfplotsset{
  thesisaxis/.style={
    width=7.2cm, height=5.6cm,
    scale only axis,
    tick align=outside,
    tick label style={font=\footnotesize},
    label style={font=\small},
    legend style={font=\footnotesize, draw=none, fill opacity=0.85,
      text opacity=1, row sep=1pt, inner sep=2pt},
    every axis plot/.append style={line width=0.9pt, mark size=1.6pt},
    grid=both,
    grid style={line width=0.2pt, draw=gray!18},
    major grid style={line width=0.3pt, draw=gray!30},
  },
  % beta 4本のカーブを一括設定する系列スタイル
  betaplotA/.style={color=betaA, mark=*,        mark options={fill=betaA}},
  betaplotB/.style={color=betaB, mark=square*,  mark options={fill=betaB}},
  betaplotC/.style={color=betaC, mark=triangle*,mark options={fill=betaC}},
  betaplotD/.style={color=betaD, mark=diamond*, mark options={fill=betaD}},
}

%% delta 軸（対数）の共通設定
\pgfplotsset{
  deltaxaxis/.style={
    xmode=log, log basis x=10,
    xtick={0.3333,1,2,3,5,10},
    xticklabels={$\tfrac13$,$1$,$2$,$3$,$5$,$10$},
    xlabel={$\frac{|\Delta \alpha|E}{p}$},
    xmin=0.28, xmax=12,
  },
}

%% beta 凡例（4系列に共通で使う）
\newcommand{\betalegend}{\legend{$\beta pE=0.5$,$\beta pE=1$,$\beta pE=2$,$\beta pE=3$}}

%% 4本の beta カーブ（誤差棒つき）を一括描画。
%%   #1 データファイル, #2 x列名, #3 値の接頭辞, #4 誤差の接頭辞
\newcommand{\plotbetaserr}[4]{%
  \addplot[betaplotA, error bars/.cd, y dir=both, y explicit]
    table[x=#2, y=#3_b05, y error=#4_b05]{#1};
  \addplot[betaplotB, error bars/.cd, y dir=both, y explicit]
    table[x=#2, y=#3_b1,  y error=#4_b1]{#1};
  \addplot[betaplotC, error bars/.cd, y dir=both, y explicit]
    table[x=#2, y=#3_b2,  y error=#4_b2]{#1};
  \addplot[betaplotD, error bars/.cd, y dir=both, y explicit]
    table[x=#2, y=#3_b3,  y error=#4_b3]{#1};
}

%% 4本の beta カーブ（誤差棒なし）。#1 ファイル, #2 x列, #3 値の接頭辞
\newcommand{\plotbetas}[3]{%
  \addplot[betaplotA] table[x=#2, y=#3_b05]{#1};
  \addplot[betaplotB] table[x=#2, y=#3_b1]{#1};
  \addplot[betaplotC] table[x=#2, y=#3_b2]{#1};
  \addplot[betaplotD] table[x=#2, y=#3_b3]{#1};
}

%% を呼び、各図は \input{analysis/fig_xxx} で取り込む。
%%
%% ※ documentclass のオプションに [dvipdfmx] が付いているので、
%%   tikz/pgfplots は自動で dvipdfmx ドライバを使う（platex + dvipdfmx 構成）。

\usepackage{tikz}
\usepackage{pgfplots}
\usepackage{pgfplotstable}
\pgfplotsset{compat=1.18}
\usepgfplotslibrary{groupplots}
\usetikzlibrary{pgfplots.groupplots}

%% データファイルの場所（本文からの相対パス）。必要なら上書き可。
\providecommand{\figdata}{analysis/data}

%% 図が多い節で float が後ろの節へ流れ込まないように配置を緩める。
\renewcommand{\topfraction}{0.9}
\renewcommand{\bottomfraction}{0.9}
\renewcommand{\textfraction}{0.08}
\renewcommand{\floatpagefraction}{0.75}
\setcounter{topnumber}{2}
\setcounter{bottomnumber}{1}
\setcounter{totalnumber}{4}

%% '#' 始まりの行をコメントとして無視させる（gnuplot 風 .dat 用）。
\pgfplotstableset{col sep=space, comment chars={\#}}

%% beta = 0.5,1,2,3 を表す色（色覚多様性に配慮した4色）
\definecolor{betaA}{HTML}{0072B2} % beta=0.5  青
\definecolor{betaB}{HTML}{009E73} % beta=1    緑
\definecolor{betaC}{HTML}{E69F00} % beta=2    橙
\definecolor{betaD}{HTML}{D55E00} % beta=3    朱

%% 図の基本スタイル
\pgfplotsset{
  thesisaxis/.style={
    width=7.2cm, height=5.6cm,
    scale only axis,
    tick align=outside,
    tick label style={font=\footnotesize},
    label style={font=\small},
    legend style={font=\footnotesize, draw=none, fill opacity=0.85,
      text opacity=1, row sep=1pt, inner sep=2pt},
    every axis plot/.append style={line width=0.9pt, mark size=1.6pt},
    grid=both,
    grid style={line width=0.2pt, draw=gray!18},
    major grid style={line width=0.3pt, draw=gray!30},
  },
  % beta 4本のカーブを一括設定する系列スタイル
  betaplotA/.style={color=betaA, mark=*,        mark options={fill=betaA}},
  betaplotB/.style={color=betaB, mark=square*,  mark options={fill=betaB}},
  betaplotC/.style={color=betaC, mark=triangle*,mark options={fill=betaC}},
  betaplotD/.style={color=betaD, mark=diamond*, mark options={fill=betaD}},
}

%% delta 軸（対数）の共通設定
\pgfplotsset{
  deltaxaxis/.style={
    xmode=log, log basis x=10,
    xtick={0.3333,1,2,3,5,10},
    xticklabels={$\tfrac13$,$1$,$2$,$3$,$5$,$10$},
    xlabel={$\frac{|\Delta \alpha|E}{p}$},
    xmin=0.28, xmax=12,
  },
}

%% beta 凡例（4系列に共通で使う）
\newcommand{\betalegend}{\legend{$\beta pE=0.5$,$\beta pE=1$,$\beta pE=2$,$\beta pE=3$}}

%% 4本の beta カーブ（誤差棒つき）を一括描画。
%%   #1 データファイル, #2 x列名, #3 値の接頭辞, #4 誤差の接頭辞
\newcommand{\plotbetaserr}[4]{%
  \addplot[betaplotA, error bars/.cd, y dir=both, y explicit]
    table[x=#2, y=#3_b05, y error=#4_b05]{#1};
  \addplot[betaplotB, error bars/.cd, y dir=both, y explicit]
    table[x=#2, y=#3_b1,  y error=#4_b1]{#1};
  \addplot[betaplotC, error bars/.cd, y dir=both, y explicit]
    table[x=#2, y=#3_b2,  y error=#4_b2]{#1};
  \addplot[betaplotD, error bars/.cd, y dir=both, y explicit]
    table[x=#2, y=#3_b3,  y error=#4_b3]{#1};
}

%% 4本の beta カーブ（誤差棒なし）。#1 ファイル, #2 x列, #3 値の接頭辞
\newcommand{\plotbetas}[3]{%
  \addplot[betaplotA] table[x=#2, y=#3_b05]{#1};
  \addplot[betaplotB] table[x=#2, y=#3_b1]{#1};
  \addplot[betaplotC] table[x=#2, y=#3_b2]{#1};
  \addplot[betaplotD] table[x=#2, y=#3_b3]{#1};
}
